\scanpage{50}
\mspage{6}

\uline{Remarques} \quad En fait, on a un foncteur plus précis, vers une catégorie de graphes munis d'une structure où la partie $\rho$ de la pondération est \deleted{\unsure{complétée}} \deleted{\unsure{remplacée}} \inserted{remplacée} par une structure plus précise, à savoir \deleted{une application} \inserted{une \deleted{\unsure{appli}} famille}
\[
 \deleted{\unsure{A}\;(\widetilde R,\alpha,p:R\to\underset{\substack{\shortparallel\\\vec A/\sigma}}{\vec A})}\quad (R_a)_{a\in A}
\]
d'ensembles, et pour tt $a\in A$, si $\omega(a)$ désigne l'ens. des \deleted{où $\alpha$ est une involution sans points fixes de} \deleted{des} ``orientations'' de $a$ (card $\omega(a)\in\{1,2\}$), une application

% slanted
\margin{$\omega(a)\overset{\mathrm{def}}{=}\{\vec a\in\vec A\,/\,\vec a\ \text{au dessus de }a\}$}

\deleted{$R$, et $p$ une application $R\to A$ comp. avec $\alpha$}
\[
 \omega\longmapsto u_\omega:\quad \omega(a)\longrightarrow\mathrm{Circ}(R_a)
 \quad\text{ens des permutations circulaires sur $R_a$}
\]
\deleted{et induisant \unsure{une iso.} une application \uline{injective}}
\[
 \deleted{R/\alpha\lhook\joinrel\longrightarrow\vec A}
\]
avec la condition que si $\omega,\omega'\in\omega(a)$, $\omega\ne\omega'$, alors $\deleted{u_\omega u_{\omega'}=\mathrm{id}_{R_a}}\quad u_{\omega'}=u_\omega^{-1}$

On peut aussi considérer que la donnée supplémentaire revient à la donnée de
\[
 \begin{aligned}
 r:&\quad A\longrightarrow\mathbb N\\
 \shortparallel&\\[-4pt]
 a&\longmapsto\mathrm{Card}\,p^{-1}(a)
 \end{aligned}
\]
(la ``pondération-arêtes''), et la donnée d'une famille de torseurs
\[
 (R_a)_{a\in A},\ \deleted{\unsure{r(a)>0}}
\]
$R_a$ un torseur sous $\mathbb Z/r(a)\mathbb Z\mathbin{\underset{\{\pm1\}}{\wedge}}^{\omega(a)}$. On associe à $(X,\widetilde\Sigma,R)$ une telle structure supplémentaire où
\[
 R_a=R\cap\widetilde\Sigma_a
 \qquad\text{($\widetilde\Sigma_a$ la comp. connexe de $\widetilde\Sigma$ \deleted{\unsure{au point}} qui correspond à $a\in A\simeq\pi_0(\widetilde\Sigma)$)}
\]
\deleted{$R=R$ (sic), $R\longrightarrow\underset{\substack{\shortparallel\\\pi_0(\Sigma)}}{S}$ déduit de $R\subset\Sigma$}

$\omega\mapsto u_\omega$ défini de façon \unsure{adéquate|évidente}.

et \deleted{$p$ déduit de l'ordre circulaire sur les $R_a$}

\uline{NB} Même si $r(\omega)=1$ ou $\unsure{r(\omega)|\text{même}}=2$, le fait de disposer
% archive sheet number: 49

